\scanpage{3}
I\quad Construction d'un $2$-polyèdre triangulé régulier orienté \ldots\hfill p.~\uline{1}

$\overline{\underline{\mathrm{II}}}$\quad Relation avec l'isom $\mathrm{Gl}(2,\mathbb{Z})\simeq\mathbb{G}_{\infty,3}$\hfill p.~$12'$

$\overline{\underline{\mathrm{III}}}$\quad Cas d'une structure $\frac12$ symplectique variable, sur $M$ : \unsure{recherche|recherches} des structures sur $P^{\deleted{\unsure{\infty|\sim}}}_{\underline M}=(S_M,F_M=\coprod_{i\in\pi}F_{M,i})$.\hfill p.~18

$\overline{\underline{\mathrm{IV}}}$\quad Récapitulation (1)\hfill p.~30

$\overline{\underline{\mathrm{V}}}$\quad Re-récapitulation, avec programme de travail\hfill p.~35

$\overline{\underline{\mathrm{VI}}}$\quad Calcul du commutant de $\mathrm{Sl}_{\Lambda}(\unsure{M|\underline M})$ dans $M^*/\pm1$, et applications\hfill p.~46

$\overline{\underline{\mathrm{VII}}}$\quad Axiomatique des structures triangulaires congruentielles (1)\hfill p.~$51'$

$\overline{\underline{\mathrm{VIII}}}$\quad Axiomatique : reformulation\hfill p.~65

$\overline{\underline{\mathrm{IX}}}$\quad Point de vue graphes à \deleted{\unsure{opérateurs|opérations}, \unsure{de graphes|de pointements}} \inserted{homothéties et multigraphes} à \deleted{\unsure{opérateurs|opérations}} \inserted{homothéties}. \inserted{``Mise en \unsure{équations|graphes}'' \unsure{gruppenthéorétique|groupe-théorétique}\ldots}\hfill p.~81.

$\overline{\underline{\mathrm{X}}}$\quad ``Antipodismes'' en général -- dualité (pour les \deleted{\unsure{ou}} graphes) -- \uline{prégraphes à homothéties}.\hfill p.~92

1. Antipodismes généralisés\hfill \deleted{p.~\unsure{111}} \deleted{\illegible{1}} 92

2. Où les étoiles prennent le pas -- structures de contour \unsure{invariante|invariantes}\dotfill 98 \deleted{\unsure{102}}

3. Présentation symétrique $S,S^\vee$ sur une $\mathbb{Z}/n\mathbb{Z}$-droite affine\dotfill 102

\deleted{4. Retour sur condition d'isomorphisme\hfill \unsure{113}}

$\overline{\underline{\mathrm{XI}}}$\quad \uline{Récapitulation} des notions obtenues\hfill 114

1. Graphes, prégraphes, \unsure{jougs|je-ges|jeux|pages}, multigraphes, \unsure{multizygons|multizygues|multipages} --\hfill 114\qquad --

2. Condition d'isomorphie pour un multigraphe\hfill $120'$

3. Condition d'homogénéité : Traductions \unsure{gruppenthéorétiques|groupe-théorétiques}\hfill \unsure{122|123}

4. Cas des graphes \unsure{triangulaires|triangulants|triangulatoires}.\hfill 13\deleted{\unsure{7}}8

\margin{p.~79/90/95\unsure{$'$|$1$}}
% archive sheet number: 2
